\section{Explicit equations for the same projective scheme}
\label{sp:sec:trinomials}
This section replaces a homogeneous circuit by a polynomial list of
three-term homogeneous equations without changing its projective scheme.
The number of equations is part of the input size and will be essential.
Projective feasibility reductions already occur in
\citet[Lemma~1 and Theorem~6]{sp-gkp}; their question asks for a nontrivial
common zero over an algebraic closure. The construction here identifies
all affine charts as rings, so it preserves nilpotents, coherent
cohomology, and rational points over every extension of the ground field.

\begin{theorem}[A projective presentation by trinomials]
\label{sp:thm:trinomials}
Let $f$ be a positive-degree homogeneous arithmetic circuit over $\Fp$ in
$r$ variables. Addition gates have inputs of equal degree, multiplication
gates are binary, and multiplication by a field element is unary.
There is a deterministic algorithm, polynomial in the circuit bit length
and $\log p$, producing homogeneous equations with at most three monomials
each in ordinary projective space whose common zero scheme is isomorphic
over $\Fp$ to $V(f)\subset\PP^{r-1}$.

After subdividing edges, let $s$ count all gates and let $h\geq1$ be their
maximum level, with every input at level zero and every edge joining
successive levels. One may use $r^2+2rs$ coordinates and at most
\[
 r^3+r^2+2r^2s+3rs+r
\]
equations. Every equation has degree at most $p^h$, and every exponent
has $O(h\log p+1)$ bits. In particular, fixed characteristic and
logarithmic circuit depth give polynomial numerical degrees.
\end{theorem}

\begin{proof}
Evaluate all degree-zero subcircuits and absorb their field values into
unary scalar gates. The remaining inputs are variables and every
remaining gate has positive homogeneous degree.
Subdividing an edge by identity gates preserves homogeneity and increases
the circuit size at most quadratically. Put $D=p^h$. If $g$ is a gate,
write $e_g$ for its homogeneous degree, $\ell_g$ for its level, and
$t_g=p^{h-\ell_g}$. Since $e_g\leq2^{\ell_g}$, we have $t_ge_g\leq D$.
Consider the following degree-$D$ forms in the original variables:
\begin{align*}
 T_j&=x_j^D,&B_{ji}&=x_j^{D-1}x_i,\\
 V_{jg}&=x_j^{D-e_g}f_g,&
 U_{jg}&=x_j^{D-t_ge_g}f_g^{t_g}.
\end{align*}
Here $B_{jj}$ means $T_j$, and $f_g$ is the polynomial computed at $g$.
The $T_j$ have no common projective zero, so these forms define a morphism
from $\PP^{r-1}$ to one ordinary projective space. Introduce coordinates
with their displayed names and impose, for every indicated index,
\begin{align}
 B_{ji}T_k^{D-1}&=B_{kj}^{D-1}B_{ki},&
 T_iT_k^{D-1}&=B_{ki}^{D},\label{sp:eq:cross1}\\
 V_{jg}T_k^{D-e_g}&=B_{kj}^{D-e_g}V_{kg},&
 U_{jg}T_k^{D-t_ge_g}&=B_{kj}^{D-t_ge_g}U_{kg},\label{sp:eq:cross2}\\
 V_{jg}^{t_g}&=T_j^{t_g-1}U_{jg}.&&\label{sp:eq:powers}
\end{align}
For a fixed anchor $j$, abbreviate its coordinates by $T,V_g,U_g$.
Impose the following two relations at each gate:
\[
\begin{array}{c|c|c}
f_g&\text{ordinary relation}&\text{power relation}\\ \hline
x_i&V_g=B_{ji}&U_g=T_i\\
f_a+f_b&V_g=V_a+V_b&U_g^p=T^{p-1}(U_a+U_b)\\
f_af_b&TV_g=V_aV_b&U_g^p=T^{p-2}U_aU_b\\
c f_a&V_g=cV_a&U_g^p=cT^{p-1}U_a
\end{array}
\]
The last row includes identity gates. These equations follow from
$t_a=t_b=pt_g$ on the preceding level and $c^{p^v}=c$ in $\Fp$.
All exponents are nonnegative, including $p-2$ when $p=2$.

We first prove that these relations alone define $\PP^{r-1}$ as a scheme.
A homogeneous prime containing all $T_j$ contains every $B_{ji}$ by
\eqref{sp:eq:cross1}. At input gates it contains every $U_{jg}$.
Induction through the circuit, using the power relations and primality,
then puts every $U_{jg}$ in the prime. At a multiplication gate in
characteristic two the right side is $U_aU_b$, so this induction still
works. Equation~\eqref{sp:eq:powers} then puts every $V_{jg}$ in the prime.
Thus no projective point has all $T_j=0$. This argument is about
homogeneous primes, and establishes an open cover of the entire scheme,
including any possible nonreduced structure.

On $D_+(T_k)$ set $T_k=1$ and $b_i=B_{ki}$; then $b_k=1$.
Equation~\eqref{sp:eq:cross1} forces
\[
 T_i=b_i^D,\qquad B_{ji}=b_j^{D-1}b_i.
\]
The ordinary gate relations at $k$ give $V_{kg}=f_g(b)$ by induction.
Equation~\eqref{sp:eq:powers} gives
$U_{kg}=f_g(b)^{t_g}$. The cross relations then give
\[
 V_{jg}=b_j^{D-e_g}f_g(b),\qquad
 U_{jg}=b_j^{D-t_ge_g}f_g(b)^{t_g}.
\]
Conversely, these substitutions satisfy every equation as polynomial
identities in $\Fp[b_i:i\ne k]$. The chart coordinate ring is therefore
exactly this polynomial ring. In particular it has no extra infinitesimal
structure. On overlaps the transition is $b_i\mapsto b_i/b_j$, the usual
transition of projective space.

Finally impose $V_{j,\mathrm{out}}=0$ for every anchor $j$.
On the $k$th chart these equations generate precisely $(f(b))$: the
$k$th output equation is $f(b)=0$, and the other equations are its
monomial multiples. The chart identifications consequently restrict to
an isomorphism with $V(f)$, with its full scheme structure.

The two cross families contribute $r^3+r^2$ and $2r^2s$ equations;
the power identifications and gate equations contribute $3rs$, and the
output contributes $r$. The displayed degree and bit bounds follow from
$e_g,t_g e_g\leq D$. Coordinates can be indexed in binary, and each
relation requires only polynomially many bits. This also proves the
algorithmic assertions.
\end{proof}

The products in the Boolean encoding, and the pure powers in
Theorem~\ref{thm:stabilizer}, have balanced homogeneous circuits of
logarithmic depth. Theorem~\ref{sp:thm:trinomials} therefore gives
polynomially many explicit trinomials of polynomial numerical degree for
the same smooth projective variety. The larger embedding has the same
two nonzero coherent groups, both of dimension one. It does not have
bounded codimension.

\begin{example}[A positive unit clause over $\mathbb F_3$]
\label{sp:ex:unit}
Use the degree-ten smooth hypersurface of Section~\ref{sec:unit-clause},
with scalar one on $H^8(\OO_X)$ and the displayed Laurent generator.
The finite count there can be separated into its three slope cases.
For rational points, $F=(yA-a)z+y^2+\sum_i x_i^2$ over $\mathbb F_3$.
Among the $3^9$ choices of $(x,y)$, the slope in $z$ is nonzero in
$9356$ cases; slope and constant term both vanish in $3517$ cases;
and only the slope vanishes in $6810$ cases. The affine cone therefore
has $9356+3\cdot3517=19907$ points. Removing its origin and dividing
by two gives
\[
 \#X(\mathbb F_3)=9953\equiv2=1+1\pmod3.
\]
The projective presentation of Theorem~\ref{sp:thm:trinomials} has exactly
the same count because it is an isomorphism over $\mathbb F_3$.
\end{example}

\section{What a bound on the number of sparse equations permits}
\label{sp:sec:bounded}
A fixed number of explicitly sparse equations makes modular point
counting elementary even when their degrees and ambient dimension grow.
The relevant operation is summing monomials on the affine cone.

\begin{theorem}[Bounded sparse systems]
\label{sp:thm:bounded}
Fix a prime $p$ and an integer $c\geq0$. The number modulo $p$ of common
rational points of at most $c$ homogeneous sparse equations in
$\PP^{r-1}_{\Fp}$ is computable in deterministic polynomial bit time.
Exponents may be binary, and no smoothness, reducedness, dimension, or
complete-intersection hypothesis is required. If $T_i$ is the term count
of equation $i$, a bound on the number of enumerated terms is
\[
 \prod_i(1+T_i^{p-1}).
\]
For trinomials the same conclusion holds for $O(\log L)$ equations,
where $L$ is the total input bit length.
\end{theorem}

\begin{proof}
Discard zero equations. A nonzero constant gives the empty scheme.
For the remaining forms $f_1,\ldots,f_b$, the polynomial
\[
 D(x)=\prod_{i=1}^b(1-f_i(x)^{p-1})
\]
is the zero-set indicator on $\Fp^r$. For a nonnegative exponent $e$,
\[
 \sum_{a\in\Fp}a^e=
 \begin{cases}-1,&e>0\text{ and }p-1\mid e,\\0,&\text{otherwise}.
 \end{cases}
\]
The $e=0$ case is zero because the field has characteristic $p$.
Consequently a term $c_\nu x^\nu$ contributes $(-1)^r c_\nu$ precisely
when every $\nu_j$ is a positive multiple of $p-1$. If $A$ is the cone
count and $B$ the projective count, then
\[
 A\equiv(-1)^r\!
 \sum_{\nu\in((p-1)\mathbb Z_{>0})^r}c_\nu\pmod p,
 \qquad A=1+(p-1)B,
 \qquad B\equiv1-A\pmod p.
\]
This calculation depends only on rational points, so it applies to any
common zero scheme.

For each factor choose either $1$ or an ordered $(p-1)$-tuple of its
monomials, with one overall minus sign. Coinciding exponent vectors may
be processed separately. A coordinate exponent is a sum of at most
$c(p-1)$ input exponents. If $E$ is the largest exponent, each sum has
$O(\log(2+c(p-1)E))$ bits. Coefficients are multiplied modulo $p$.
Thus every choice takes polynomial time in $L$. If the coordinates are
indexed sparsely and $r$ itself is binary, merge the indices occurring in
the chosen terms and first check whether every one of the $r$ coordinates
occurs; this avoids a loop of length $r$ when it is unnecessary.

Ordered choices can be grouped by multiplicity. There are
$\binom{T_i+p-2}{p-1}$ multiplicity vectors in $f_i^{p-1}$.
For a trinomial this is $\binom{p+1}{2}$, a constant at fixed $p$.
For $b\leq C\log L$ the total is at most
$(1+\binom{p+1}{2})^{C\log L}$, which is polynomial in $L$.
\end{proof}

The prime-field monomial summation is the precision-one form of the
modular algorithms of \citet[\S3.1, pp.~5--6]{sp-ggl}.
Their \S3.3, pp.~7--8, replaces an extension-field equation by
prime-field component equations and combines their zero indicators by
$1-\prod_i(1-f_i^{p-1})$. Here the same product keeps the dependence on
the number of equations explicit. Wan's sparse hypersurface theorem,
discussed above, is consistent with this bound.

\begin{corollary}[An obstruction to sparse compression]
\label{sp:cor:compression}
Assume $\mathrm{Mod}_p\mathrm P\not\subseteq\mathrm P$ for a fixed prime
$p$. There is no deterministic polynomial-time transformation of every
formula into a polynomial-length explicit sparse projective system with
a fixed number of equations whose rational point count determines its
satisfying-assignment count modulo $p$ by the congruence in
Theorem~\ref{thm:main}. The same obstruction holds with
$O(\log L)$ trinomials. In particular, it applies to scheme-isomorphic
presentations of the rank-one hypersurfaces.
\end{corollary}
\begin{proof}
Such a transformation followed by Theorem~\ref{sp:thm:bounded} would
decide whether the satisfying-assignment count is nonzero modulo $p$
in polynomial time. The decision problem is
$\mathrm{Mod}_p\mathrm P$-complete, so this contradicts the stated
assumption.
\end{proof}

The conclusion concerns a uniform construction and its input length.
It is not an unconditional lower bound on the smallest number of terms
in a particular equation, nor does it rule out a long sparse list in
bounded codimension. A different projective variety with the same
coherent scalar falls under the argument when its scalar is recoverable
from its point count by a specified polynomial-time formula. One cannot
replace that last hypothesis by an unspecified coherent group.

\section{Reembedding a smooth hypersurface}
\label{sp:sec:reembedding}
The Picard group restricts how the same smooth hypersurface can acquire
a different projective embedding. It separates a change of linear
coordinates from a genuine change of polarization.

\begin{theorem}[Polarization and ambient dimension]
\label{sp:thm:reembedding}
Let $X^n\subset\PP^{n+1}_k$ be a smooth hypersurface over an algebraically
closed field, with $n\geq3$ and degree $\delta\geq2$. Write
$H=\OO_X(1)$. For every closed immersion $j:X\hookrightarrow\PP^M$,
there is an integer $a\geq1$ with $j^*\OO(1)\simeq H^{\otimes a}$.
If $a=1$, the map is the original embedding followed by a linear
inclusion and a projective coordinate change. If $a\geq2$, then
\[
 M\geq2n+1,\qquad \deg j(X)=a^n\delta.
\]
For each $a\geq2$ an embedding in $\PP^{2n+1}_k$ exists. Every set of
homogeneous equations defining such a nonlinear reembedding, even only
set-theoretically, therefore has at least $n+1$ members.
\end{theorem}

\begin{proof}
The projective complete-intersection theorem in
\citet[Expos\'e~XII, Corollaire~3.7, p.~121]{sp-sga2} gives
$\operatorname{Pic}(X)=\mathbb Z[H]$ for dimension at least three.
The theorem is stated for projective complete intersections more
generally; smoothness is more than it requires. Its proof applies
Corollaire~3.6 inductively, using the vanishing of $H^1(X,H^v)$ and
$H^2(X,H^v)$ in that dimension, together with the parafactoriality
results of Expos\'e~XI, Th\'eor\`eme~3.13. The dimension hypothesis
cannot be lowered to two: a smooth quadric surface has Picard group
$\mathbb Z^2$. Positivity gives $a>0$.

The restriction sequence for degree-one forms shows
$h^0(X,H)=n+2$. If the sections defining an embedding with $a=1$ spanned
a proper subsystem, its image would span at most $\PP^n$. An integral
closed $n$-dimensional subscheme of $\PP^n$ is all of $\PP^n$, of degree
one, whereas $H^n=\delta$. Thus the full system is used. For all $a$,
the degree is $(aH)^n=a^n\delta$.

For $a\geq2$, put $V_i=H^0(X,H^{\otimes i})$ and embed $X$ by its
complete system in $\PP(V_a^*)=\PP^R$. If $x,y$ are distinct points,
choose trivializations and denote their evaluations on $V_i$ by
$e_{x,i},e_{y,i}$. An interior point of their secant is
$u=\lambda e_{x,a}+\mu e_{y,a}$ with $\lambda\mu\ne0$.
Multiplication of sections gives
\[
 C_u:V_1\longrightarrow V_{a-1}^*,\qquad
 C_u(s)(t)=u(st).
\]
As $H^{a-1}$ is very ample, $e_{x,a-1}$ and $e_{y,a-1}$ are independent.
Therefore
\[
 \ker C_u=\{s\in V_1:s(x)=s(y)=0\}.
\]
This kernel recovers the original line $\ell=\langle x,y\rangle$ in
$\PP^{n+1}$. For a general pair, $\ell$ is not contained in $X$:
choose $y$ outside the tangent hyperplane at $x$, possible because $X$
spans $\PP^{n+1}$. The intersection $X\cap\ell$ then has finite length
$\delta$. Every decomposition of $u$ as a secant of two distinct points
uses a pair from this finite set; each pair determines the coefficients
up to common scale. The incidence of ordered distinct pairs and an
interior secant point has dimension $2n+1$ and finite general fibres
over its image. Hence the closed secant variety has dimension $2n+1$.
No separability or characteristic-zero tangent-space formula enters
this argument.

The closed secant variety also contains exactly the spans of all
length-two subschemes relevant to projection. Indeed
$\operatorname{Hilb}^2(X)$ is projective, and the universal length-two
family gives a rank-two bundle of values of $H^a$. Its projective bundle
maps properly to $\PP(V_a^*)$ with image the union of these spans.
A nonreduced length-two subscheme is a tangent direction. Smooth local
coordinates give a smooth curve through that direction, on which the
double point is a limit of two distinct points. Thus this proper image
is the secant closure just computed.

An arbitrary subsystem $W\subset V_a$ gives projection from
$\PP(W^\perp)$. It defines a closed immersion precisely when its sections
separate every length-two subscheme, equivalently when the center misses
the closed secant variety. This is the point-and-tangent criterion
\citep[Lemmas~33.23.1--2, tag~0E8R]{sp-stacks}.
If $M'=\dim W-1$, the projective dimension theorem now gives
$M'\geq2n+1$. If the center is empty, the same inequality follows from
$R\geq\dim\operatorname{Sec}(X)=2n+1$. Conversely a general linear
center of dimension $R-2n-2$ misses the secant variety, so projection
to $\PP^{2n+1}$ is an embedding. Finally the codimension there is
$n+1$, and the height theorem bounds the number of defining equations
below by the codimension.
\end{proof}

For $a=2$, \citet[Proposition~1.3(a), Theorem~3.3, and Example~6.6]{sp-zak}
gives the corresponding quadratic secant computation: positive original
codimension forces first secant defect zero. The contraction above
also treats higher powers and arbitrary characteristic directly.
A matrix argument identifying quadratic forms with sums of squares
requires care in characteristic two; multiplication of sections does
not. The algebraic-closure existence of a general center does not
construct a rational center over the original finite field. Quantitative
field and output issues are addressed in Section~\ref{sp:sec:compression}.

\section{Anticanonical complete intersections and ideal generation}
\label{sp:sec:ideal}
For a global complete intersection the coherent coefficient is still a
single coefficient, and sparse generators of its homogeneous ideal can
be used without first choosing a regular sequence.

Let $R=\Fp[x_0,\ldots,x_N]$, $\mathfrak m=(x_0,\ldots,x_N)$, and let
$X\subset\PP^N$ be a smooth geometrically integral global complete
intersection of codimension $c$ and dimension $n=N-c\geq1$.
Choose a homogeneous regular sequence $g_1,\ldots,g_c$ for $I_X$,
with degrees $e_i>0$ satisfying
\[
 \sum_{i=1}^c e_i=N+1.
\]
This is the anticanonical degree condition. The ideal generated by a
regular sequence is saturated here: its quotient is Cohen--Macaulay
of dimension $n+1\geq2$, so it has no $\mathfrak m$-torsion.
For a $c$-tuple $u$ whose degrees sum to $N+1$, set
\[
 C(u)=[(x_0\cdots x_N)^{p-1}](u_1\cdots u_c)^{p-1}.
\]

\begin{proposition}[The coherent coefficient and point count]
\label{sp:prop:ci-coefficient}
The structure-sheaf cohomology of $X$ is $\Fp$ in degrees $0,n$ and
zero in intermediate degrees. With its Koszul residue generator,
Frobenius on $H^n(X,\OO_X)$ has scalar $\alpha=C(g)$, and
\[
 \#X(\Fp)\equiv1+(-1)^n\alpha\pmod p.
\]
\end{proposition}
\begin{proof}
The Koszul resolution has terms $\OO(-\sum_{i\in I}e_i)$.
Every nonempty proper subset has degree sum at most $N$, so all of
its cohomology vanishes. The last term is $\OO(-N-1)$, with its single
top cohomology class $1/(x_0\cdots x_N)$. The resulting connecting
isomorphism identifies it with the top coherent line of $X$.
The map from the Koszul resolution of $(g_1^p,\ldots,g_c^p)$ to that
of $(g_1,\ldots,g_c)$ multiplies its top term by
$(g_1\cdots g_c)^{p-1}$. Applying $p$th power to the inverse monomial
and then this multiplier gives the asserted coefficient.
This is the rank-one specialization of
\citet[Proposition~4.2.1, pp.~18--19, and Algorithm~II, \S4.3]{sp-kudo}.
In the general regular-sequence formula, negative exponent vectors
$k^{(i)}$ of sum $-\sum e_j$ index a basis, and the $(j,i)$ entry is
the coefficient with exponent $-p k^{(i)}+k^{(j)}$ in
$(\prod g_j)^{p-1}$.

For the point count, expand the cone indicator
$\prod_i(1-g_i^{p-1})$. Every surviving monomial under summation on
$\Fp^{N+1}$ has each exponent a positive multiple of $p-1$ and therefore
total degree at least $(N+1)(p-1)$. Only the full product reaches this
degree; its surviving exponent vector is exactly $(p-1,\ldots,p-1)$.
The cone count is consequently $(-1)^{c+N+1}\alpha$ modulo $p$.
Since the cone count is $1+(p-1)\#X(\Fp)$, the projective count is
$1+(-1)^{N-c}\alpha$ modulo $p$.
\end{proof}

\begin{theorem}[Sparse generators before saturation]
\label{sp:thm:generators}
Suppose nonzero homogeneous sparse forms $f_1,\ldots,f_s$ generate
$I_X$ as an ideal, and $c$ is known. Enumerate the $c$-element sublists
$u$ whose degrees sum to $N+1$. If any $C(u)$ is nonzero, return that
value; if all vanish, return zero. This computes $\alpha$.
For fixed $p,c$, its bit cost is
\[
 T^{c(p-1)}\operatorname{poly}(L),
\]
where $T$ is the total input term count and $L$ its bit length.
\end{theorem}
\begin{proof}
The images of $g_1,\ldots,g_c$ form a homogeneous basis of
$I_X/\mathfrak m I_X$. They span by definition; if a constant linear
combination vanished in that quotient, graded Nakayama would remove a
generator, contrary to height $c$
\citep[Lemma~10.56.1(4), tag~0EKB]{sp-stacks}.
The classes of the input forms also
span, so some $c$-element sublist is a basis. Its degree multiset is
$\{e_1,\ldots,e_c\}$, hence it is enumerated.

For any enumerated tuple write $u=Ag$, with $A_{ij}$ homogeneous of
degree $\deg u_i-e_j$, and zero if that number is negative. Every
nonzero determinant term has degree
$\sum_i\deg u_i-\sum_j e_j=0$, so $\det A\in\Fp$.
The transformation law needed here is
\begin{equation}
 \prod_i(AY)_i^{p-1}
   =(\det A)^{p-1}\prod_jY_j^{p-1}
 \quad\text{in }R[Y_1,\ldots,Y_c]/(Y_1^p,\ldots,Y_c^p).
 \label{sp:eq:determinant-law}
\end{equation}
Indeed, the top graded piece of this truncated polynomial ring is the
line spanned by $\prod Y_j^{p-1}$. Over the fraction field of $R$,
elementary shears and permutations fix that line with multiplier one,
while a diagonal substitution multiplies it by the determinant to
power $p-1$. Gaussian elimination proves the assertion for invertible
matrices; a singular matrix reduces to a diagonal matrix with a zero
entry and gives zero. The free monomial basis descends the identity
to $R$.

Substitute $Y_j=g_j$. Every error term is divisible by some $g_j^p$,
which belongs to $(x_0^p,\ldots,x_N^p)$ because $g_j$ has positive
degree. Such a term cannot contribute to the target coefficient.
Therefore
\[
 C(u)=(\det A)^{p-1}\alpha.
\]
Every nonzero determinant in $\Fp$ has $(p-1)$st power one. For the
basis sublist, the reduction of $A$ modulo $\mathfrak m$ is invertible,
so $C(u)=\alpha$. If $\alpha=0$, all candidates vanish; otherwise a
nonzero candidate exists and every such candidate has the right value.

If form $i$ has $t_i$ terms, enumeration of its $(p-1)$st power takes
$t_i^{p-1}$ ordered choices. Summing over $c$-element sublists gives at
most
\[
 \left(\sum_i t_i^{p-1}\right)^c\leq T^{c(p-1)}
\]
choices. Degree tests, coefficient operations, and addition of binary
exponents each have polynomial bit cost.
\end{proof}

\begin{example}[Why saturation changes the coefficient search]
\label{sp:ex:cubic}
Over $\mathbb F_5$ consider
\[
 g=y^2z-x^3-xz^2.
\]
On $z=1$, the cubic $x^3+x=x(x-2)(x-3)$ has distinct roots; since the
characteristic is not two, this proves affine smoothness. Its unique
point at infinity is $[0:1:0]$, where $g_z=1$.
In $g^4$, the coefficient of $x^4y^4z^4$ uses two copies of $y^2z$ and
one copy of each negative monomial. Thus
\[
 \alpha=\frac{4!}{2!1!1!}=12\equiv2\pmod5.
\]
For $x=0,1,2,3,4$, the values of $x^3+x$ are $0,2,0,0,3$.
Only the three zeros are squares, and each gives one affine point.
The point at infinity gives $\#X(\mathbb F_5)=4\equiv1-2$.

The quartics $xg,yg,zg$ generate $\mathfrak m(g)$, whose saturation is
$(g)$. They define exactly the same projective scheme, yet there is no
cubic sublist, so the algorithm in Theorem~\ref{sp:thm:generators} would
return zero if its ideal-generation hypothesis were omitted. Keeping
only $yg$ adds the line $y=0$: it has six rational points, three already
on the cubic, so that zero set has seven points. The next theorem
recovers the scalar from all three quartics and treats a larger class
of presentations requiring saturation.
\end{example}

\section{Recovering scalars and cofactor tuples from sparse multiples}
\label{sp:sec:saturation}
The failure of the preceding sublist algorithm does not make every
presentation requiring saturation hard. For hypersurfaces, a bound on
the degrees of the unknown cofactors gives an explicit recovery algorithm.

\begin{lemma}[Removing monomial content and Frobenius powers]
\label{sp:lem:preprocessing}
Suppose a nonzero homogeneous sparse list over $\Fp$ has saturation
$(g)$, where $g$ is geometrically integral of degree greater than one.
One may replace any input by its quotient by its largest monomial
divisor. One may also replace an input all of whose exponent vectors
are divisible by $p$ by its polynomial $p$th root.

More generally, fix $e\geq1$ and write each exponent vector uniquely
as $r+p^e q$, with $0\leq r_k<p^e$. Regard the coefficients of one
input as a sparse matrix indexed by $r$ and $q$. If that matrix has
rank one, reconstruct
\[
 f_i=U(x)V(x)^{p^e}
\]
and replace $f_i$ by $U(x)V(x)$ whenever $V$ is nonconstant.
Repeatedly applying any available operations preserves the saturated
ideal, never increases term counts, and terminates in polynomial time
in the sparse bit length. The rank-one and whole-power replacements
also preserve any radical saturated ideal generated by the list after
saturation.
\end{lemma}
\begin{proof}
No coordinate divides $g$. Thus the monomial content of $f_i=g h_i$
belongs entirely to $h_i$, and the divided polynomial still belongs
to $(g)$. Over $\Fp$, an input with all exponents divisible by $p$
is a $p$th power. Primality of $(g)$ implies that its $p$th root
also belongs to $(g)$.

For the general operation put $P=p^e$. A nonzero rank-one coefficient
matrix has Cartesian support: if its nonzero row and column sets have
sizes $u$ and $v$, then all $uv$ entries in their product are nonzero.
Thus $uv$ is the original term count. Choose a nonzero pivot, factor
the matrix as a column times a row, and use these coefficients as those
of $U$ and $V$. Coefficients lie in $\Fp$, so their $P$th powers are
themselves, giving $f_i=UV^P$ exactly. Homogeneity and the Cartesian
support imply that all row exponents have one common total degree
and all column exponents have one common total degree. Hence $U,V$
are homogeneous. Expanding $UV$ uses only $uv$ term pairs, and
collecting them cannot increase the original term count.

If $I$ is a radical saturated ideal and $UV^P\in I$, then
$(UV)^P=U^{P-1}(UV^P)\in I$ implies $UV\in I$.
Moreover $UV^P$ is a multiple of $UV$. In every operation the old
ideal $J$ and its enlargement $J'$ therefore satisfy
\[
 J\subset J'\subset I,
\]
where $I=(g)$ for monomial stripping. Since $J^{\mathrm{sat}}=I$,
the same equality holds for $J'$.

We give an iteration bound because the exponents may be binary.
Let $n$ be the number of coordinate indices actually occurring in
the input being processed, and put $a=\deg U$ and $b=\deg V\geq1$.
No operation introduces a new coordinate index, and $n\leq L$ even
when the ambient dimension is recorded sparsely. Residue exponents
are smaller than $P$, so $a\leq n(P-1)$. If the degree before a
rank-one step is $D=a+Pb$, its decrease is $(P-1)b$, and
\[
 D\leq n(P-1)b+2(P-1)b=(n+2)(P-1)b.
\]
Each nontrivial step therefore decreases the degree by at least the
fraction $1/(n+2)$. There are $O(n\log(D+1))$ such steps from initial
degree $D$. Whole-power extraction is included in this bound, and
complete monomial stripping is needed at most once between such steps.
For each step one tests only
$e\leq\lfloor\log_p(\max\hbox{ exponent})\rfloor$; larger $e$
give a constant $V$. The sparse rank-one test checks Cartesian support
and the pivot factorization in polynomial time. All exponents decrease
and their bit lengths remain bounded by the input bit length.
This proves the claimed total complexity.
\end{proof}

The degree hypotheses in the following recovery results may be checked
after this preprocessing. Thus arbitrarily large binary monomial
multipliers and repeated Frobenius powers do not by themselves violate
those hypotheses.

\begin{theorem}[Explicit recovery with bounded total cofactor degrees]
\label{sp:thm:bounded-cofactors}
Fix $p$ and $B\geq0$. Let $g\in\Fp[x_0,\ldots,x_N]$ have degree
$\delta=N+1$ and define a smooth geometrically integral hypersurface of
positive dimension. Suppose a nonempty list of nonzero homogeneous
sparse forms satisfies
\[
 (f_1,\ldots,f_s):\mathfrak m^\infty=(g),\qquad
 \delta\leq\deg f_i\leq\delta+B.
\]
From this list alone, the tuple of cofactors $f_i/g$ can be recovered
up to a common nonzero scalar, and the coherent Frobenius scalar can be
computed, in deterministic polynomial bit time. The same conclusion holds when the
degree bounds hold after Lemma~\ref{sp:lem:preprocessing}. The algorithm does not expand $g$.
More generally, its running time is polynomial in $L$ and
\[
 U=\sum_i\binom{N+b_i}{b_i},\qquad b_i=\deg f_i-\delta,
\]
followed by at most $t_j^{p-1}\binom{N+b_j}{b_j}$ coefficient choices
for one selected input form $f_j$.
\end{theorem}
\begin{proof}
Since each $f_i$ belongs to $(g)$, write $f_i=g h_i$, where $h_i$ is
homogeneous of degree $b_i$. These cofactors are unknown. Put
$K=(h_1,\ldots,h_s)$. As $g$ belongs to the saturation of $gK$, there
is an $a$ with $\mathfrak m^a g\subset gK$. Cancellation in the
polynomial domain gives $\mathfrak m^a\subset K$. Thus the $h_i$
have no common projective zero and have greatest common divisor one.

Introduce an unknown coefficient for every monomial of degree $b_i$
in a form $H_i$. Solve the homogeneous linear equations
\begin{equation}
 f_1 H_i-f_i H_1=0\qquad(2\leq i\leq s).
 \label{sp:eq:cofactor-system}
\end{equation}
There are $U$ unknowns. Coefficients of the products can be collected
from sparse lists with polynomial work in $L,U$; no factorization or
expansion of $g$ is used. The vector $(h_i)$ is a nonzero solution.
We claim the solution space is a line. In the fraction field, every
solution is $(H_i)=q(h_i)$. Write $q=a/b$ with coprime polynomials.
Since all $H_i$ are polynomials, $b$ divides every $h_i$ and is therefore
constant. Hence $q$ is a polynomial. The prescribed homogeneous degrees
force $q$ to be a scalar. Consequently any nonzero solution is
$H_i=\lambda h_i$ with $\lambda\in\Fp^*$.

The absence of a common projective zero implies that some $H_j$ is
nonzero at $e_0=[1:0:\cdots:0]$. Choose such a $j$ by reading the
coefficient $c$ of $x_0^{b_j}$. Dehomogenize at $x_0=1$ and write bars
for the resulting polynomials. Then $\bar f_j=\bar g'\bar H_j$,
where $g'=g/\lambda$.
In the finite ring
\[
 A=\Fp[z_1,\ldots,z_N]/(z_1^p,\ldots,z_N^p)
\]
we have $\bar H_j^p=c^p$ because all its nonconstant monomials vanish
after taking $p$th powers. Therefore
\begin{equation}
 (\bar g')^{p-1}
   =c^{-p}\bar f_j^{p-1}\bar H_j\quad\text{in }A.
 \label{sp:eq:truncated-quotient}
\end{equation}
The coefficient of $(z_1\cdots z_N)^{p-1}$ on the left is the desired
homogeneous coefficient: homogeneity of $g'$ forces the missing
$x_0$ exponent to be $p-1$. Scaling by $\lambda$ does not change it
over $\Fp$, since $\lambda^{p-1}=1$.

Compute the right side by choosing an ordered $(p-1)$-tuple of terms
of $f_j$ and one term of $H_j$. Accept a choice precisely when the
exponent sum in every coordinate $x_1,\ldots,x_N$ is $p-1$; the
$x_0$ exponent is ignored after dehomogenization. Multiply the sum by
$c^{-p}$. Equation~\eqref{sp:eq:truncated-quotient} proves correctness.
Since $B$ is fixed, $U\leq s\binom{N+B}{B}$ is polynomial in the input
length, and so are the linear system and coefficient enumeration.
This remains true with sparsely indexed coordinates: smoothness and
$\deg g=N+1>1$ force $g$ to involve every coordinate, since an omitted
coordinate would give a singular vertex. Degrees in each coordinate
add under multiplication, so every nonzero $f_i=g h_i$ also involves
every coordinate. Hence $N+1$ is bounded by the displayed input length.
\end{proof}

For the cubic in Example~\ref{sp:ex:cubic}, the cofactor system recovers
$(x,y,z)$ up to a common scalar. Choosing the cofactor $x$ and setting
$x=1$ reduces the calculation to the coefficient of $y^4z^4$ in
$g(1,y,z)^4$, again giving $2$. Thus generation only after saturation
is tractable in this bounded-excess case. Unbounded cofactor degrees
and general fixed codimension remain distinct questions: the linear
system above then has no polynomial bound on its monomial universe,
and a complete intersection with several generators need not have
scalar cofactors at all.

\subsection{Recovery from one multiple of small degree excess}
The preceding linear system bounds every cofactor degree. A univariate
restriction argument needs this bound for only one listed multiple.
The saturation hypothesis is stronger than a greatest-common-divisor
condition: after dehomogenization, the cofactors generate the unit ideal,
and this remains true on every affine line.

\begin{lemma}[Evaluation of a cofactor by univariate gcds]
\label{sp:lem:line-gcd}
Suppose $F_i=G H_i\in\Fp[z_1,\ldots,z_N]$ and
$(H_1,\ldots,H_s)=(1)$. Assume that a selected nonzero explicit sparse input
$F_j$ has total degree at most $D$. Over a finite extension $K/\Fp$
with $|K|>D$, one can choose $a\in K^N$ with $F_j(a)\ne0$ and
construct an exact evaluation algorithm for
\[
 H_j(z)/H_j(a).
\]
Each evaluation has bit cost polynomial in $D$, the sparse bit length
of all $F_i$, and the explicit field representation. No numerical bound
on the degrees of the other inputs is required.
\end{lemma}
\begin{proof}
Find $a$ by successively specializing the variables of the known sparse
polynomial $F_j$. At each stage the current polynomial is nonzero and
has degree at most $D$ in the next variable. At most $D$ choices can
make it the zero polynomial in the remaining variables: choose any
nonzero coefficient polynomial in that variable to see this bound.
Since $|K|>D$, testing any $D+1$ distinct field values and collecting
equal sparse terms finds a specialization preserving nonzeroness.
This takes polynomial work in $D$ and the explicit field representation.
For the stated applications one may choose the least power of $p$
exceeding $D$, so that $|K|\leq p(D+1)$.

For a query $z\in K^N$, put $v=z-a$ and
\[
 q_i(t)=F_i(a+tv),\qquad G_v(t)=G(a+tv),\qquad
 H_{i,v}(t)=H_i(a+tv).
\]
The chosen $q_j$ is nonzero, has degree at most $D$, and has
$q_j(0)\ne0$. The unit-ideal identity for the $H_i$ restricts to one
for the $H_{i,v}$, so
\[
 \gcd_i q_i(t)=G_v(t)
\]
up to a nonzero constant. This holds for every query line, including
lines on which degrees drop. Moreover $G_v$ is nonzero because
$G_v(0)H_{j,v}(0)=q_j(0)\ne0$.

Construct $q_j$ as a dense univariate polynomial of degree at most
$D$. For every other input, compute $r_i=q_i\bmod q_j$ directly in
$K[t]/(q_j)$. A sparse monomial becomes a product of powers of the
linear polynomials $a_k+t v_k$; repeated squaring uses a number of
operations polynomial in the binary exponent length, with all
remainders of degree less than $D$. Thus one never constructs the
possibly enormous $q_i$. Compute the monic polynomial
\[
 d_v=\gcd(q_j,r_1,\ldots,r_s).
\]
There is a nonzero constant $\lambda_v$ with
$G_v=\lambda_v d_v$, and hence
\[
 Q_v=q_j/d_v=\lambda_v H_{j,v},\qquad
 \frac{Q_v(1)}{Q_v(0)}=\frac{H_j(z)}{H_j(a)}.
\]
The denominator is nonzero. If $q_j$ is constant, the same formula
holds with monic gcd one. Univariate arithmetic and Euclidean gcds
satisfy the claimed bound.
\end{proof}

This use of univariate gcds is related to classical black-box polynomial
algorithms. In particular, \citet[Section~3, pp.~310--311]{sp-kaltofen-trager}
construct general gcd evaluation procedures by randomized projections
and address the possibility of extra factors appearing after
specialization. Here the unit-ideal identity ensures correct gcds on
every affine line. The preceding proof supplies the fixed-characteristic
algorithm and its sparse-bit bound directly.

\begin{theorem}[One multiple of bounded degree excess]
\label{sp:thm:one-cofactor}
Fix $p$ and $B\geq0$. Let a nonempty homogeneous sparse list over
$\Fp$ have saturation $(g)$, where $g$ is a smooth geometrically
integral anticanonical hypersurface equation in $\PP^N$ of positive
dimension. If at least one listed nonzero form satisfies
\[
 \deg f_j\leq N+1+B,
\]
then the coherent Frobenius scalar is computable in deterministic
polynomial bit time. The other input degrees may be arbitrarily large
and binary-encoded. The same assertion holds when this degree bound
is obtained after Lemma~\ref{sp:lem:preprocessing}.
\end{theorem}
\begin{proof}
Write $f_i=g h_i$. As in
Theorem~\ref{sp:thm:bounded-cofactors}, saturation implies
$\mathfrak m^a\subset(h_i)$ for some $a$. Set $x_0=1$ and denote the
resulting polynomials by $F_i,G,H_i$. Then $(H_i)=(1)$.
The selected cofactor has homogeneous degree
$b_j=\deg f_j-(N+1)\leq B$; its dehomogenization has total degree
at most $B$.

Choose the least power $q$ of $p$ exceeding $N+1+B$ and work over
an explicit $K=\mathbb F_q$. Its extension degree is $O(\log(N+B+2))$,
and even exhaustive search for an irreducible defining polynomial
has polynomial cost at fixed $p,B$. Lemma~\ref{sp:lem:line-gcd}
finds $a$ with $F_j(a)\ne0$ and evaluates $H_j/H_j(a)$ at any point.

We recall a deterministic interpolation procedure to keep the degree
bound explicit. Choose $B+1$ distinct elements
$u_0,\ldots,u_B\in K$. For every multi-index $\alpha\in\mathbb N^N$
with $|\alpha|\leq B$, query the polynomial at
$(u_{\alpha_1},\ldots,u_{\alpha_N})$.
Use the Newton basis
\[
 P_\beta(z)=\prod_{i=1}^N\prod_{k=0}^{\beta_i-1}(z_i-u_k),
 \qquad |\beta|\leq B.
\]
Ordered by total degree, its evaluation matrix on these points is
triangular: $P_\beta(u_\alpha)=0$ if some $\beta_i>\alpha_i$,
and its diagonal entries are nonzero. There are
$\binom{N+B}{B}$ basis elements. Solve this system and expand the
basis elements, each of which has at most $2^B$ monomials.
Homogenize to the known degree $b_j$. Dividing by one nonzero
coefficient produces a homogeneous $H\in\Fp[x_0,\ldots,x_N]$:
indeed it is the true $h_j$ divided by the corresponding original
$\Fp$ coefficient. Thus $f_j=g'H$, where $g'$ differs from $g$ by
a nonzero $\Fp$ scalar.

It remains to compute the scalar without expanding a translated input.
Put $c=H(1,a)\ne0$. In the ring
$K[z_1,\ldots,z_N]/(z_1^p,\ldots,z_N^p)$,
\[
 g'(1,a+z)^{p-1}
   =c^{-p}F_j(a+z)^{p-1}H(1,a+z).
\]
The homogeneous coordinate transformation
$x_i\mapsto x_i+a_i x_0$ has determinant one and preserves the Hasse
coefficient. To verify this directly, the top-degree line of
$K[x_0,\ldots,x_N]/(x_0^p,\ldots,x_N^p)$ is spanned by
$(x_0\cdots x_N)^{p-1}$, and a linear transformation $A$ acts on it
by $\det(A)^{p-1}$. This identity holds for diagonal matrices,
permutations and elementary shears, which generate the general linear
group. It therefore holds over $K$, and our shear has factor one.
Homogeneity of $g'$ then identifies the desired scalar with the
coefficient of $(z_1\cdots z_N)^{p-1}$ in the displayed identity.

Choose an ordered $(p-1)$-tuple of terms of $f_j$ and one term of
$H$, and let $e_i$ be the resulting exponent sum in $x_i$ for
$i\geq1$. After the affine translation its contribution to the target,
in addition to its coefficient, is
\[
 \prod_{i=1}^N\binom{e_i}{p-1}a_i^{\,e_i-p+1},
\]
with contribution zero if any $e_i<p-1$. Sum these contributions and
multiply by $c^{-p}$. This uses at most
$t_j^{p-1}\binom{N+B}{B}$ choices. It neither expands $g'$ nor the
translated $f_j$. The final scalar belongs to $\Fp$, and replacing
$g$ by $g'$ leaves it unchanged because their scaling ratio has
$(p-1)$st power one. All steps have polynomial sparse-bit cost.
\end{proof}

We spell out the sparse-interpolation ingredient. It is the
Klivans--Spielman method \citep[Section~6.3, Theorem~11]{sp-ks},
with a variable scaling parameter as in
\citet[Section~3, Theorem~3.2]{sp-cs-sparse}. The variable parameter
makes the characteristic-independent reconstruction explicit.

\begin{lemma}[Sparse interpolation with numerical degree bounds]
\label{sp:lem:sparse-interpolation}
Let $P\in K[z_1,\ldots,z_n]$ have total degree at most $b$ and at
most $R$ nonzero terms. Given exact evaluations over a sufficiently
large explicit extension of $K$, its sparse list can be recovered
deterministically using $\operatorname{poly}(n,b,R)$ evaluations and
field operations. The required field cardinality is polynomial in
$n,b,R$. This assertion holds in every characteristic.
\end{lemma}
\begin{proof}
The zero-variable or degree-zero cases are immediate. Assume
$n,b,R\geq1$, put $m=(n-1)\binom R2+1$, and choose an ordinary
integer prime $\ell>\max(b,m)$ of polynomial size. For
$t=0,\ldots,m-1$ define integer weights
\[
 w_k(t)=t^{k-1}\bmod\ell\in\{0,\ldots,\ell-1\}
 \qquad(1\leq k\leq n).
\]
For two different exponent vectors $\alpha,\beta$ of total degree
at most $b$, the polynomial
\[
 \sum_{k=1}^n(\alpha_k-\beta_k)T^{k-1}\in\mathbb F_\ell[T]
\]
is nonzero, since $\ell>b$, and has at most $n-1$ roots. Fewer than
$m$ values of $t$ can therefore cause a collision for any pair of
monomials in the support of $P$. For at least one tested $t$, the
integer weights $\sum_k\alpha_k w_k(t)$ are distinct on that support.
The auxiliary prime $\ell$ is unrelated to the characteristic of $K$.

Each univariate image
$P(y^{w_1(t)},\ldots,y^{w_n(t)})$ has degree at most
$D=b(\ell-1)$. Interpolate all these images over a field with more
than $D$ elements and choose an image with the largest number of
nonzero terms. A substitution cannot increase the term count, and
one tested substitution separates the whole support. Thus a maximizing
choice is itself collision-free, even when other choices exhibit
coefficient cancellation. If every image is zero, then $P=0$.

Fix a maximizing weight vector $w$. For each coordinate $k$,
interpolate the bivariate polynomial
\[
 P(y^{w_1},\ldots,y^{w_{k-1}},u y^{w_k},
                 y^{w_{k+1}},\ldots,y^{w_n}).
\]
Its degrees in $y,u$ are at most $D,b$. For every distinct label
$e=\sum_i\alpha_iw_i$, the coefficient of $y^e$ is exactly
$c_\alpha u^{\alpha_k}$. It reveals $\alpha_k$ and $c_\alpha$.
Reading these coefficients for all $k$ recovers every exponent vector.
Ordinary univariate and bivariate interpolation uses at most
$m(D+1)+n(D+1)(b+1)$ evaluations, a polynomial bound.
The proof uses no fixed scalar whose powers might collide in small
characteristic. At fixed $p$, a common extension of $\Fp$ with the
required polynomial cardinality can be constructed deterministically
in polynomial time by searching its polynomially many monic candidate
defining polynomials and testing irreducibility.
\end{proof}

\begin{theorem}[A designated sparse cofactor]
\label{sp:thm:sparse-cofactor}
Fix $p$. In the principal saturation problem of
Theorem~\ref{sp:thm:one-cofactor}, designate an index $j$ with $f_j\ne0$ and supply a
valid bound $R$ on the number of monomials of $h_j=f_j/g$.
Put $b=\deg f_j-(N+1)$. The scalar is computable deterministically
in $\operatorname{poly}_p(L,N+b,R)$ bit operations, even when the
other input degrees are arbitrary binary integers. The algorithm also
recovers $h_j$ up to a nonzero $\Fp$ scalar.

This is a polynomial sparse-bit bound when $b$ and $R$ are
polynomially bounded in $L$. It does not assert polynomial dependence
on $\log b$ or $\log R$. The index and support promise must refer
to the same selected input. All conditions may be imposed after
Lemma~\ref{sp:lem:preprocessing}.
\end{theorem}
\begin{proof}
The case $b=0$ is immediate from
Theorem~\ref{sp:thm:one-cofactor}. Otherwise, dehomogenization leaves
$H_j$ with at most $R$ terms and individual degrees at most $b$.
Lemma~\ref{sp:lem:line-gcd} constructs an exact evaluation algorithm
for $H_j/H_j(a)$ in time polynomial in $L,N+b$ for each query.
Instead of total-degree interpolation, apply
Lemma~\ref{sp:lem:sparse-interpolation}. Choose its common extension
also to satisfy $|K|>N+1+b$ before constructing the evaluation oracle.
At fixed $p$, construction and arithmetic have polynomial cost in
$N+b,R$. The line evaluation algorithm supplies every required
evaluation, including at zeros of $G$ or $H_j$.

The recovered polynomial is exactly $H_j/H_j(a)$ under the supplied
valid sparsity bound. Homogenize it to degree $b$ and normalize one
nonzero coefficient, obtaining a cofactor $\widehat h_j$ over $\Fp$ as in
Theorem~\ref{sp:thm:one-cofactor}. Write $f_j=g'\widehat h_j$, where
$g'$ differs from $g$ by a nonzero $\Fp$ scalar.

For the scalar calculation put $\mathbf 1=(1,\ldots,1)$ and define
the $\Fp$-linear Cartier operator
on monomials by
\[
 \mathcal C_p(x^\nu)=
 \begin{cases}
 x^{(\nu-(p-1)\mathbf 1)/p},&
          \nu_i\equiv p-1\pmod p\text{ for every }i,\\
 0,&\text{otherwise}.
 \end{cases}
\]
The monomial definition gives
$\mathcal C_p(AH^p)=H\mathcal C_p(A)$ for polynomials over $\Fp$.
Since $\deg g'=N+1$, homogeneity and Katz's coefficient formula give
$\mathcal C_p((g')^{p-1})=\alpha$, the coherent scalar. Hence
\[
 \mathcal C_p(f_j^{p-1}\widehat h_j)
   =\mathcal C_p((g')^{p-1}\widehat h_j^p)
   =\alpha\widehat h_j.
\]
Compute the left side from at most $t_j^{p-1}R$ term choices and
read its coefficient at the monomial normalized to coefficient one
in $\widehat h_j$. This gives $\alpha$ without expanding $g'$ and
proves the stated running time.
\end{proof}

\begin{corollary}[An unspecified sparse cofactor]
\label{sp:cor:unspecified-cofactor}
Fix $p$ and let the principal saturation presentation satisfy the
geometric hypotheses of Theorem~\ref{sp:thm:one-cofactor}. Suppose
that some nonzero residual cofactor, whose index is not supplied, has total
degree at most $b$ and at most $R$ nonzero terms. Given these bounds
and $0<\epsilon<1$, there is a Monte Carlo algorithm returning the
coherent scalar correctly with probability at least $1-\epsilon$.
Every run has bit cost
$\operatorname{poly}_p(L,N+b,R,\log(1/\epsilon))$.
The other cofactor degrees are unrestricted binary integers.
The bounds may be imposed after Lemma~\ref{sp:lem:preprocessing}.
\end{corollary}
\begin{proof}
Discard zero input rows. Try every remaining row whose displayed degree satisfies
$0\leq b_i:=\deg f_i-(N+1)\leq b$. There is a promised good row
among them. If $b=0$, any such row gives the deterministic computation
of Theorem~\ref{sp:thm:one-cofactor}, so assume $b\geq1$.
Choose a common explicit extension $K/\Fp$ large enough for
Lemma~\ref{sp:lem:sparse-interpolation}, with $|K|>N+1+b$ and
$|K|>4b$. Its cardinality is polynomial in $N+b,R$.
For each eligible $i$, Lemma~\ref{sp:lem:line-gcd} constructs the
exact evaluation oracle for
$P_i(z)=h_i(1,z)/h_i(1,a_i)$, where $f_i(1,a_i)\ne0$.
The polynomial $P_i$ has total degree at most $b_i$, even if it has
more than $R$ terms.

Run the explicit interpolation procedure with parameters $b_i,R$.
Its fixed finite lists of univariate and bivariate evaluations have
the stated polynomial bound without a sparsity promise. Reject a
row if an image has more than $R$ terms, if a labelled bivariate
coefficient is not a single monomial with consistent coefficient, or
if the resulting candidate $C_i$ has more than $R$ terms or degree
greater than $b_i$. Also require $C_i(a_i)=1$. A promised good row
returns its exact $P_i$ and passes every such check.

For a candidate, compare its evaluations with the oracle at
\[
 k=\left\lceil\log_2(s/\epsilon)\right\rceil
\]
fresh independent random points, where $s$ is the number of input
rows. If $C_i\ne P_i$, their difference is a nonzero polynomial of
degree at most $b$. For uniform coordinates its probability of
vanishing is at most $b/|K|$. A bounded-bit implementation also
suffices: identify the field elements with $0,\ldots,|K|-1$, choose
$\lceil\log_2|K|\rceil$ random bits, and reduce the resulting integer
modulo $|K|$ to select an element. Each value has probability at most
$2/|K|$. For independent coordinates with maximum point mass $\mu$,
the root bound is $b\mu$: induct on the variables, separating failure
of the leading coefficient from the roots of each remaining univariate
specialization. It therefore gives vanishing probability
at most $2b/|K|<1/2$. Thus an incorrect candidate passes all comparisons
with probability at most $2^{-k}$, and a union bound over the rows
bounds any false acceptance by $\epsilon$.

For an accepted candidate, homogenize to degree $b_i$, normalize a
nonzero coefficient, and require that every normalized coefficient
belongs to $\Fp$. The condition at $a_i$ ensures that the translated
quotient calculation has a nonzero denominator. Compute the scalar
as in Theorem~\ref{sp:thm:one-cofactor}, reject the row if this value
does not belong to $\Fp$, and otherwise return it. The good row is always
accepted. Outside the false-acceptance event, the first accepted row
is its true normalized cofactor, so the returned scalar is correct.
All queries, checks, and random-bit counts satisfy the displayed bound.
\end{proof}

\begin{corollary}[Explicit recovery under degree and support bounds]
\label{sp:cor:sparse-tuple}
If valid degree and support bounds $b,R$ apply to every cofactor,
the entire cofactor tuple is recoverable, up to one common nonzero
$\Fp$ scalar, in $\operatorname{poly}_p(L,N+b,R)$ bit time.
\end{corollary}
\begin{proof}
Discard zero input forms, whose cofactors are zero, and renumber the
remaining forms. Apply Theorem~\ref{sp:thm:sparse-cofactor} to every
remaining index. Its outputs
are $H_i=\lambda_i h_i$ with $\lambda_i\in\Fp^*$.
For each $i$, the sparse products $f_1H_i$ and $f_iH_1$ are nonzero
constant multiples. Reading and checking their coefficient ratio
allows one to rescale $H_i$ so that $f_1H_i=f_iH_1$.
All components then have the same normalization.
\end{proof}

\subsection{Coordinate degrees and access to hidden cofactor ratios}
A bound on individual degrees of the displayed multiples is quite different
from a bound on individual degrees of the hidden cofactors.

\begin{lemma}[Coordinate degrees forced by smoothness]
\label{sp:lem:coordinate-degree}
Let a homogeneous polynomial $g$ of degree $\delta>1$ define a smooth
hypersurface in $\PP^N$ over a field. Then
\[
 \delta-1\leq\deg_{x_j}g\leq\delta\qquad(0\leq j\leq N).
\]
For every nonzero polynomial $h$, one has
$\deg_{x_j}(gh)=\deg_{x_j}g+\deg_{x_j}h$.
Consequently a fixed individual-degree bound $D$ on even one nonzero
multiple of a smooth anticanonical equation forces $N\leq D$.
If all the cofactors of a principal saturation presentation have
individual degrees at most $b$, then for each coordinate point some
cofactor has total degree at most $b$ and is nonzero at that point.
\end{lemma}
\begin{proof}
At $e_j$, either the coefficient of $x_j^\delta$ is nonzero, or
$e_j$ lies on the hypersurface. In the second case smoothness requires
some nonzero first derivative there, which can only be supplied by a
monomial $x_j^{\delta-1}x_k$, $k\ne j$. This proves the lower bound.
The degree identity follows by treating the polynomials as univariate
in $x_j$ over the domain generated by the other variables. For the
last assertion, the cofactors have no common projective zero. A
homogeneous cofactor nonzero at $e_j$ contains a pure power
$x_j^{\deg h_i}$; its exponent is at most $b$.
\end{proof}

\begin{corollary}[Polynomial scalar computation from bounded individual degrees]
\label{sp:cor:individual-scalar}
Fix $p$ and $b$. For the smooth principal saturation problem, if every
cofactor has individual degrees at most $b$, then the coherent scalar
is computable in deterministic polynomial bit time. The same holds
if, after Lemma~\ref{sp:lem:preprocessing}, every input has individual
degrees at most $N+b$.
\end{corollary}
\begin{proof}
Some cofactor is nonzero at $e_0$, so
Lemma~\ref{sp:lem:coordinate-degree} bounds its total degree by $b$.
Theorem~\ref{sp:thm:one-cofactor} applies to the corresponding input.
The stated input bound implies the cofactor bound by the same lemma.
\end{proof}

The last observation does not say that the small-total-degree cofactors
have no common projective zero. Thus it does not by itself reduce
bounded individual cofactor degree to the linear-system proof in
Theorem~\ref{sp:thm:bounded-cofactors}; the one-cofactor theorem is
what gives the stronger corollary.
There is nevertheless exact access to their rational ratios, even at
points where the common hypersurface factor vanishes.

\begin{lemma}[Rational cofactor oracle]
\label{sp:lem:cofactor-oracle}
Let $f_i=g h_i$ be a homogeneous principal saturation presentation as
in Theorem~\ref{sp:thm:bounded-cofactors}, with no bound on cofactor
degrees. Dehomogenize at $x_0=1$, writing $F_i,G,H_i$ for the resulting
polynomials. From the sparse input one can simulate, over any explicit
finite extension of $\Fp$, simultaneous evaluation of
\[
 H_1/H_j,\ldots,H_s/H_j
\]
at any query point, returning a pole symbol exactly when $H_j=0$.
Each query has polynomial bit cost in the input length and extension
representation. One can also choose $j$ for which $H_j(0)\ne0$.
\end{lemma}
\begin{proof}
The $H_i$ generate the unit ideal: a power of the irrelevant ideal
lies in $(h_i)$, and dehomogenizing its power of $x_0$ gives $1$.
If some $F_i(a)$ is nonzero, then $G(a)\ne0$. In that case
$H_j(a)\ne0$ precisely when $F_j(a)\ne0$, and the ordinary ratios
$F_i(a)/F_j(a)$ give the required answers.

If all $F_i(a)$ vanish, then $G(a)=0$, since the $H_i(a)$ do not all
vanish. Smoothness of the affine hypersurface gives
$\nabla G(a)\ne0$, and differentiation gives
\[
 \nabla F_i(a)=H_i(a)\nabla G(a).
\]
Hence $H_j(a)=0$ exactly when $\nabla F_j(a)=0$. Otherwise choose
$k$ with $\partial_kF_j(a)\ne0$ and return
$\partial_kF_i(a)/\partial_kF_j(a)$ for each $i$.
This is also a deterministic way to select a denominator nonzero at
$a=0$: choose a nonzero $F_j(0)$ if one exists, and otherwise choose a
nonzero gradient. Such a $j$ exists by the preceding argument.
Sparse polynomials and their first derivatives are evaluated by binary
exponentiation, so binary exponents cause no additional degree factor.
\end{proof}

In a rational-interpolation application the simultaneous tuple matters:
$\gcd(H_1,\ldots,H_s)=1$, while separately reducing individual ratios
can change the denominator and enlarge its support. Once a simultaneous
interpolation procedure recovers the $H_i$ up to a common scalar,
normalize $H_j(0)=1$. Uniqueness then implies that the normalized
coefficients belong to $\Fp$, even if interpolation used an extension.
Equation~\eqref{sp:eq:truncated-quotient} computes the scalar with at
most $t_j^{p-1}\#\operatorname{supp}(H_j)$ term choices.
The oracle remains efficient when the numerical cofactor degrees are
large, whereas the line-gcd algorithm is polynomial in a selected
numerical degree. A reconstruction theorem for these rational ratios
with polynomial dependence on binary degree lengths could therefore
address a different regime. For polynomially bounded degrees and
supports, Corollary~\ref{sp:cor:sparse-tuple} already recovers the
whole tuple using ordinary sparse interpolation. Characteristic-zero
rational-interpolation results alone do not supply a further
fixed-characteristic conclusion.

\begin{theorem}[Explicit recovery of all bounded-individual-degree cofactors]
\label{sp:thm:individual-cofactors}
Fix a prime $p$ and an integer $b\geq0$. Suppose a nonempty sparse
homogeneous list over $\Fp$ satisfies
\[
 (f_1,\ldots,f_s):\mathfrak m^\infty=(g),\qquad
 \deg g=N+1,
\]
where $g$ defines a smooth geometrically integral hypersurface of positive
dimension. Write $f_i=g h_i$ and assume
$\deg_{x_j}h_i\leq b$ for every $i,j$.
The entire tuple $(h_i)$, up to one common nonzero scalar, can be
recovered as explicit sparse polynomials in
\[
 \exp\bigl(O_{p,b}((\log(L+2))^2)\bigr)
\]
deterministic bit operations, where $L$ is the sparse input length.
In particular this also computes the coherent scalar. No sparsity bound
on the unknown cofactors is assumed, and $g$ is not expanded. The
cofactor bound may be imposed after
Lemma~\ref{sp:lem:preprocessing}.
In particular, the conclusion applies when the displayed inputs satisfy
$\deg_{x_j}f_i\leq N+b$ for all $i,j$.
\end{theorem}
\begin{proof}
Discard zero input forms and return their cofactors as zero; renumber
the remaining nonempty list.
The case $b=0$ follows from
Theorem~\ref{sp:thm:bounded-cofactors}; suppose $b\geq1$.
We use the explicit candidate-list theorem of
\citet[Theorems~1.7 and~5.1; Remark~1.9]{sp-bkrs}. For a $T$-term
polynomial in $n$ variables of individual degree at most $D$, that
algorithm lists all factors of individual degree at most $b$ having
no nonconstant monomial divisor, up to constant associates. It may
include spurious candidates. The list has at most
\[
 C=\left(\sum_{k=0}^{b}\binom Dk\right)^{\lceil\log_2T\rceil}
\]
entries, each an explicit sparse polynomial with at most
$T^{4b^2\lceil\log_2n\rceil}$ terms; increasing these bounds by
polynomial factors has no effect here. Its running time is polynomial
in these quantities and $n,D$, with the stated univariate-factorization
cost. The arbitrary-field theorem, rather than the characteristic
restriction of their separate circuit-list theorem, is the one used.

Take $n=N+1$, $T=\max(2,\max_i t_i)$ and $D=N+1+b$.
These parameters bound every input individual degree, because individual
degrees add and $\deg_{x_j}g\leq N+1$. They are polynomially bounded
in $L$ by Lemma~\ref{sp:lem:coordinate-degree}.
Let $M_i$ be the largest monomial dividing $f_i$, read directly from
its exponent list. Since $g$ is geometrically integral of degree greater
than one, it has no coordinate divisor. Thus $M_i$ is exactly the
largest monomial dividing $h_i$.
Apply the candidate-list algorithm to $f_i/M_i$, include the constant
candidate $1$, and multiply all candidates by $M_i$. Delete zero
candidates and retain only those homogeneous of degree
\[
 b_i=\deg f_i-(N+1)
\]
and with individual degrees at most $b$. The resulting list
$\mathcal L_i$ contains a nonzero constant multiple of $h_i$.

The unwanted candidates can be removed simultaneously without a
polynomial-divisibility test. Enumerate $H_1\in\mathcal L_1$.
For each $i\geq2$, search $U_i\in\mathcal L_i$ for which the two
nonzero sparse products $f_1U_i$ and $f_iH_1$ are proportional by a
constant. The proportionality is found from any nonzero coefficient
and checked by equality of the collected sparse lists. If such a
candidate exists, rescale it to obtain
\[
 f_1H_i=f_iH_1.
\]
If the search succeeds for every $i$, keep the tuple. Some tuple passes,
since the true cofactors occur in the lists. Conversely, every passing
tuple has $H_i=q h_i$ in the common fraction field. The cofactors
have greatest common divisor one by saturation, so the denominator
of $q$ is constant. The prescribed homogeneous degrees then force
$q$ itself to be constant. Hence every passing tuple is correct up
to a common scalar. There are at most $O(sC^2)$ sparse product
comparisons, with polynomial factors allowed in the list bounds.
There is no enumeration of the Cartesian product of all lists.

The published list algorithm can use a sufficiently large finite
extension. Its sparse interpolation and constant-variate steps need
an extension of cardinality bounded by a polynomial in the displayed
list and support bounds. An extension degree
$O_b((\log(L+2))^2)$ therefore suffices. Constructing it even by
exhaustive search for an irreducible polynomial costs at most
quasipolynomial time. The deterministic univariate factorization
bound over $\mathbb F_{p^e}$ is polynomial in the degree, $e$ and $p$
\citep[Theorem~2.13]{sp-bkrs}, so it fits the same bound at fixed $p$.
The comparisons above take place in this common explicit extension.

Choose a component nonzero at $e_0$ and divide the whole passing tuple
by its value there. The normalized tuple equals
$(h_i/h_j(e_0))_i$ and consequently has coefficients in $\Fp$.
Using its $j$th component, whose dehomogenization has constant term
one, equation~\eqref{sp:eq:truncated-quotient} reads the scalar by
$t_j^{p-1}\#\operatorname{supp}(H_j)$ term choices. The resulting
hypersurface equation is $h_j(e_0)g$, whose scalar is the same as
that of $g$ because $h_j(e_0)^{p-1}=1$.
All list sizes, support sizes and operations are
$\exp(O_{p,b}((\log(L+2))^2))$.
Finally, the displayed input condition implies
$\deg_{x_j}h_i\leq b$ by
Lemma~\ref{sp:lem:coordinate-degree}.
\end{proof}

\begin{remark}\label{sp:rem:individual-vs-total}
Bounded individual cofactor degree allows total cofactor degrees to grow
with the dimension. For example, the cofactors
\[
 x_1+x_0,\ldots,x_N+x_0,\quad x_0x_1\cdots x_N
\]
have individual degrees at most one and no common projective zero.
The first $N$ alone vanish at $[1:-1:\cdots:-1]$, and the last
cofactor has degree $N+1$. Thus one cannot simply omit the high-degree
cofactors and invoke the bounded-total-degree algorithm.
\end{remark}

\section{Many linear equations: a characteristic boundary}
\label{sp:sec:linears}
Linear equations can compactly represent products of linear forms inside
one sparse nonlinear equation. This has a different counting effect in
characteristic two and at every odd prime.

\begin{theorem}[Binary linear constraints]
\label{sp:thm:binary}
Fix $c,b\geq1$. Given any number of homogeneous linear equations over
$\mathbb F_2$ and at most $c$ additional homogeneous sparse equations,
one can compute the projective point count modulo $2^b$ in deterministic
polynomial bit time. Coordinates are explicit and exponents may be
binary.
\end{theorem}
\begin{proof}
Let $Av=0$ be the linear system. Discard zero equations; a nonzero
constant makes the zero set empty. In characteristic two a monomial
on a binary vector is one precisely when every coordinate in its
support is one. Thus the integer sum of a product of monomial values
over $\ker A$ is the number of solutions of $Av=0$ with a set of
coordinates pinned to one. Gaussian elimination gives either zero or
$2^{r-\rho}$, where $\rho$ is the rank of that augmented affine system
in $r$ variables.

For parity, expand $\prod_{i=1}^c(1+f_i)$ on $\ker A$ and apply this
observation to every monomial. For higher precision, write a nonlinear
equation as the parity of its monomial values $m_1,\ldots,m_T$.
Its exact integer zero indicator is
\[
 \frac{1+\prod_{i=1}^T(1-2m_i)}2
 =1+\sum_{j=1}^T(-1)^j2^{j-1}
             \sum_{|J|=j}\prod_{i\in J}m_i.
\]
Terms with $j>b$ vanish modulo $2^b$. Multiply the $c$ truncated
indicators and sum each term by the rank calculation. There are at
most $\prod_i(1+\sum_{j=1}^b\binom{T_i}{j})$ choices, polynomial at
fixed $c,b$. Finally, the origin is a cone point, and every nonzero
binary vector is its own normalized projective representative.
Subtracting one gives the projective count modulo $2^b$.
\end{proof}

\begin{proposition}[A determinant over binary extension fields]
\label{sp:prop:binary-extension}
Let $k=\mathbb F_{2^h}$ be explicitly represented. Suppose linear
equations cut out a subspace with $r$ homogeneous coordinates, and a
fixed number $c$ of sparse forms restrict there to a smooth
positive-dimensional complete intersection with degree sum $r$.
The entry of absolute, semilinear Frobenius on its top coherent line,
in a computed Koszul residue basis, can be computed in polynomial time
in the sparse input length and field representation. Its $2^h$-power
Frobenius scalar is the norm of that entry to $\mathbb F_2$, and is
nonzero exactly when the entry is nonzero.
\end{proposition}
\begin{proof}
Choose linear coordinates on the subspace by Gaussian elimination.
Every original coordinate restricts to an explicit linear form in
$r$ variables. In the product of the $c$ nonlinear equations choose
one monomial from each. Its total degree is $r$, so its restriction
is its coefficient times a product of exactly $r$ linear forms,
counted with multiplicity. The coefficient of the squarefree monomial
in this product is the permanent of the matrix of their coefficients.
In characteristic two the permanent equals the determinant. Repeated
linear factors give repeated rows and therefore contribute zero.
Sum these determinants over the polynomially many monomial choices.
Proposition~\ref{sp:prop:ci-coefficient}, valid over perfect fields
with semilinear Frobenius, identifies the result $\beta$ with the
entry in the residue basis of the reduced complete intersection.
Its action is $a\mapsto\beta a^2$. Iterating $h$ times gives
\[
 \beta\beta^2\cdots\beta^{2^{h-1}}
   =N_{k/\mathbb F_2}(\beta).
\]
A nonzero norm lies in $\mathbb F_2^*=\{1\}$; zero has norm zero.
This is an algorithm for the coherent entry over $k$. The integer
point-count algorithm of Theorem~\ref{sp:thm:binary} specifically
counts points over the binary prime field.
\end{proof}

\begin{theorem}[Odd-characteristic hardness with one nonlinear equation]
\label{sp:thm:odd-linears}
Fix an odd prime $p$. From a width-three Boolean formula $\Phi$ one
can construct a smooth geometrically integral complete intersection
\[
 X\subset\PP^{d+k-1}_{\Fp}
\]
of dimension $d-2$ and multidegree $(1^k,d)$. Its $k$ independent
linear equations have at most three terms each; its single nonlinear
equation has at most $d+p$ terms; all equations have polynomial bit
length. Frobenius on $H^{d-2}(X,\OO_X)$ has scalar
$\#\operatorname{SAT}(\Phi)\bmod p$. Deciding whether that scalar is
nonzero is $\mathrm{Mod}_p\mathrm P$-complete on this smooth promise.
\end{theorem}
\begin{proof}
Valiant's parsimonious counting construction first turns accepting
computations into width-three formulas. His permanent construction
then gives an $r\times r$ matrix $M$ with entries in
$\{-1,0,1,2,3\}$ and
\[
 \operatorname{per}(M)=4^t S,\qquad
 t=2\,\#\{\text{literal occurrences}\}
      -\#\{\text{clauses}\},
\]
where $S=\#\operatorname{SAT}(\Phi)$
\citep[Lemmas~3.1--3.2, p.~193]{sp-valiant}.
Unused declared variables receive tautological clauses, an empty clause
can be replaced by a fixed contradiction, and a forced fresh variable
ensures a nonempty instance. These changes preserve $S$.
Put
\[
 L_i=\sum_jM_{ij}x_j,\qquad Z=\prod_{j=1}^r x_j,
 \qquad P=4^{-t}\prod_{i=1}^rL_i.
\]
The scalar $4^{-t}$ exists because $p$ is odd. The squarefree coefficient
in the row product selects one column from each row and each column
once, so $[Z]P=S$ in $\Fp$.

Adjoin $p-1$ tags and form the degree-$u$ polynomial in $u=r+p-1$
variables
\[
 G=-t_1^{p-1}P+\sum_{i=2}^{p-1}t_i^{p-1}Z.
\]
The target coefficient in $G^{p-1}$ forces each tag summand once.
It is $-(p-1)![Z]P=S$ by Wilson's theorem. Apply the smooth
stabilization, choosing a power $Q$ of $p$ with $m=Q-1>u$ and
$m\geq p$, and let $d=Q+1$.

Introduce one coordinate $w_i$ for each row form $L_i$ and replace
$\prod_i L_i$ in the displayed nonlinear polynomial by the monomial
$\prod_i w_i$. Realize $w_i=L_i$ by successive partial sums of the
nonzero terms in that row. Every equation determines one new variable
from earlier ones and has at most three terms. A zero row is represented
by $w_i=0$. Thus the $k$ equations are independent and cut out a linear
$\PP^{d-1}$ scheme-theoretically. There are at most $r^2$ new coordinates.
Restriction of the nonlinear equation is exactly the smooth stabilized
hypersurface. This proves geometric smoothness, integrality, and the
complete-intersection assertion, not merely a smoothness promise for
the constructed outputs.

The tagged part contributes $p-1$ monomials. The remaining terms are
$yzx_1^m$, $y^{Q+1}$, the $m$ pure $x_i^{Q+1}$ terms, and $-x_1z^Q$.
The total is at most $(p-1)+(m+3)=d+p$.
The anticanonical degree sum is $k+d=(d+k-1)+1$.
Elimination of the linear equations identifies the coherent line with
that of the stabilized hypersurface, of scalar $S$.

Hardness follows from the parsimonious accepting-path encoding.
For membership, a verifier guesses a normalized projective vector,
checks every listed equation by repeated squaring, and counts precisely
$\#X(\Fp)$ accepting paths. Adding $p-1$ paths gives
$\#X(\Fp)-1$ modulo $p$. By
Proposition~\ref{sp:prop:ci-coefficient}, this is a nonzero constant
multiple of the coherent scalar on the promise, so it decides exactly
its nonvanishing. The independent linear codimension $k$ grows with
the instance. It is therefore outside the bounded-total-equation
hypothesis of Theorem~\ref{sp:thm:bounded}.
\end{proof}

\begin{example}[A ternary permanent]
\label{sp:ex:permanent}
Take the $2\times2$ all-ones matrix over $\mathbb F_3$.
Then $[ab](a+b)^2=2$. Impose $w_1=a+b$ and $w_2=a+b$, and put
\[
 A=(-t_1^2w_1w_2+t_2^2ab)x_5x_6x_7x_8+a^8,
 \qquad
 F=yzA+y^{10}+\sum_{i=1}^8x_i^{10}-az^9,
\]
where $x=(a,b,t_1,t_2,x_5,x_6,x_7,x_8)$.
This is a smooth complete intersection of multidegree $(1,1,10)$ in
$\PP^{11}$, with coherent scalar $2$.
For each $(x,y)$ the row equations fix $w_1,w_2$, and $F$ is
$(yA-a)z+y^2+\sum_i x_i^2$. The slope is nonzero in $9388$ cases;
both slope and offset vanish in $3581$; only the slope vanishes in
$6714$. Hence the cone count is $9388+3\cdot3581=20131$, and
\[
 \#X(\mathbb F_3)=10065\equiv0=1+2\pmod3.
\]
\end{example}

\subsection{A support bound at every prime}
The odd-prime hardness construction uses monomials involving many
coordinates. A logarithmic bound on that support restores tractability,
even though the number of linear constraints remains arbitrary.
Discard zero equations first; a nonzero constant gives the empty
scheme and is handled separately.

\begin{theorem}[Support-sensitive modular counting]
\label{sp:thm:support}
Fix $p,c$. Let homogeneous linear equations over $\Fp$ cut out a vector
space $V$ of dimension $m$, and let at most $c$ additional nonzero
homogeneous sparse equations of positive degree have at most
$\sigma$ distinct coordinates in each
monomial. Their projective zero set has point count modulo $p$
computable in
\[
 p^{c(p-1)\sigma}\operatorname{poly}(L)
\]
bit operations. If $m>c(p-1)\sigma$, its affine cone count is zero
modulo $p$ and its projective count is one modulo $p$.
In particular $\sigma=O(\log L)$ gives polynomial time at every fixed
prime. On a positive-dimensional anticanonical complete-intersection
promise, the coherent scalar is zero under the displayed strict
inequality.
\end{theorem}
\begin{proof}
Expand the nonlinear zero indicator on $V$. At fixed $p,c$ it has
polynomially many ordered monomial choices. Each selected product
involves a coordinate set $I$ of size at most $c(p-1)\sigma$.
Its value depends only on the projection
$\pi_I:V\to\Fp^I$, whose rank $\rho$ is computed by Gaussian
elimination. Every value in its image has $p^{m-\rho}$ inverse images.
If $m>\rho$, its contribution to the sum is therefore zero modulo $p$.
Otherwise $m=\rho\leq|I|$, and enumerate the at most $p^\rho$ points
of the image, evaluating each monomial by repeated squaring.
This proves the running time and, more precisely, a separate rank bound
for each selected support. If $m>c(p-1)\sigma$, every contribution
vanishes. The cone relation gives the projective residue one; the
complete-intersection trace formula then forces the coherent scalar
to vanish.
\end{proof}

This bound measures the monomials in the displayed coordinates, before
linear elimination. It does not assert that a monomial with large
support is always hard. In the permanent construction the supports grow
linearly with the number of row variables and padding coordinates;
Theorem~\ref{sp:thm:support} isolates a tractable parameter range of the
same presentation model.

The anticanonical degree condition permits a substantially different
parameter bound. It is useful even when the monomials involve a number
of coordinates proportional to the dimension.

\begin{theorem}[Counting near the minimum possible support]
\label{sp:thm:support-strip}
Fix a prime $p$ and an integer $c$. Suppose homogeneous linear equations
over $\Fp$ cut out a vector space of dimension $m$, and the restrictions
of sparse homogeneous forms $f_1,\ldots,f_c$, of degrees
$\delta_1,\ldots,\delta_c$, define a smooth positive-dimensional
complete intersection with
$\sum_i\delta_i=m$.
In each $f_i$, discard monomials having an individual ambient-coordinate
exponent at least $p$. If any resulting form is zero, the coherent
Frobenius scalar is zero. Otherwise let $\sigma_i$ be the largest
number of distinct ambient coordinates in a remaining monomial, and put
\[
 K=(p-1)\sum_{i=1}^c\sigma_i-m.
\]
If $K<0$, the scalar is zero. If $K\geq0$, it can be computed in
\[
 T^{c(p-1)}p^K\operatorname{poly}(L)
\]
bit operations, where $T$ is the total number of displayed monomials
and $L$ is the sparse input length. More precisely, the exponential
factor for an individual monomial choice can be replaced by
$p^\rho$, where $\rho$ is the rank of a matrix with at most $K$ rows.
Consequently
\[
 \sum_i\sigma_i\leq\frac{m}{p-1}+O(\log L)
\]
is a polynomial-time range at every fixed prime and fixed $c$.
In characteristic two every potentially nonzero choice has $K=0$,
recovering the determinant computation.
\end{theorem}
\begin{proof}
Choose coordinates $w_1,\ldots,w_m$ on the linear solution space.
Every original coordinate restricts to a linear form in $w$.
By Proposition~\ref{sp:prop:ci-coefficient}, the desired scalar is
the coefficient of $(w_1\cdots w_m)^{p-1}$ in
$(f_1\cdots f_c)^{p-1}$ after restriction.
There are at most $T^{c(p-1)}$ ordered monomial choices in this
expansion. We compute each choice separately, including its product
of input coefficients.

Work in
\[
 C=\Fp[w_1,\ldots,w_m]/(w_1^p,\ldots,w_m^p).
\]
For every linear form $L$, the identity
$L^p=\sum_j a_j^p w_j^p$ makes $L^p$ zero in $C$.
This proves the asserted deletion rule. It also means that, after
combining equal ambient coordinates in a chosen product, the choice
vanishes unless it has the form
\[
 \prod_{j=1}^t L_j^{e_j},\qquad
 1\leq e_j\leq p-1,\qquad
 \sum_{j=1}^t e_j=m(p-1).
\]
Here the $L_j$ are the restrictions of the $t$ distinct ambient
coordinates used by the choice; some of the resulting linear forms
may coincide. The degree equality gives $t\geq m$, while the support
bound gives
\[
 t\leq(p-1)\sum_i\sigma_i=m+K.
\]
There is therefore no surviving choice if $K<0$.

If the $L_j$ span a space of dimension less than $m$, a linear change
of variables makes the product depend on fewer than $m$ variables.
Its degree $m(p-1)$ exceeds the largest nonzero degree of the
corresponding smaller truncated polynomial ring, so it is zero in
$C$. Otherwise choose $m$ independent $L_j$, relabel them first, and
use them as coordinates $z_1,\ldots,z_m$. The coefficient on the
top degree of $C$ transforms by the $(p-1)$st power of the coordinate
determinant, as in~\eqref{sp:eq:determinant-law}; this factor is one
over $\Fp$. Put $k=t-m\leq K$ and write the remaining forms as the
rows of $Az$, with $A$ a $k\times m$ matrix. The coefficient in
question is thus the top coefficient of
\[
 H(z)=\left(\prod_{i=1}^m z_i^{e_i}\right)
       \left(\prod_{j=1}^k (Az)_j^{e_{m+j}}\right).
\]

For every homogeneous polynomial of degree $m(p-1)$ over $\Fp$,
the coordinate power-sum identities give
\[
 [(z_1\cdots z_m)^{p-1}]H
   =(-1)^m\sum_{z\in\Fp^m} H(z).
\]
Indeed, a monomial can have nonzero sum only when every exponent
is a positive multiple of $p-1$. At this total degree every exponent
must then be exactly $p-1$.

The sum can be evaluated without enumerating $\Fp^m$.
For $s\in\Fp^k$, maintain the weighted distribution
\[
 D_i(s)=
 \sum_{\substack{(z_1,\ldots,z_i)\in\Fp^i\\
                  \sum_{a=1}^i A_{\bullet a}z_a=s}}
       \prod_{a=1}^i z_a^{e_a},
 \qquad D_0(0)=1.
\]
Each new coordinate gives the update
\[
 D_i(s)=\sum_{u\in\Fp}u^{e_i}
           D_{i-1}(s-uA_{\bullet i}).
\]
Finally the required coefficient is
\[
 (-1)^m\sum_{s\in\Fp^k}D_m(s)
                    \prod_{j=1}^k s_j^{e_{m+j}}.
\]
The computation takes $O(mp^{k+1})$ field operations and $O(p^k)$
storage, with polynomial overhead for linear algebra. All states lie
in $\operatorname{im}(A)$; choosing a basis of this image replaces
$k$ by $\rho=\operatorname{rank}(A)$ in the number of states.
At fixed $p$ this proves the claimed bound. For $p=2$, the surviving
exponents are all one, so their total degree forces $t=m$ and $k=0$.
The coefficient of the corresponding product of $m$ linear forms
is its determinant, as already used in
Proposition~\ref{sp:prop:binary-extension}.
\end{proof}

The weighted distributions in this proof are ordinary convolution
computations in the group algebra $\Fp[\Fp^k]$. Related rank-parameter
dynamic programs and coefficient-extraction algorithms for linear
satisfiability are discussed by \citet[Section~4; see also the discussion
before Theorem~2]{sp-kkn}. Here only the convolution recurrence is
needed. In particular, the argument does not invoke a character
transform in characteristic $p$, where the group order is not
invertible.

The logarithmic bound on $K$ is a genuine restriction. The hard
instances can have support arbitrarily close, in relative terms, to
the minimum in Theorem~\ref{sp:thm:support-strip}.

\begin{theorem}[Hardness with small proportional support excess]
\label{sp:thm:support-hardness}
Fix an odd prime $p$ and a real constant $\varepsilon>0$.
The nonvanishing problem in Theorem~\ref{sp:thm:odd-linears} remains
$\mathrm{Mod}_p\mathrm P$-complete when every monomial of the single
nonlinear degree-$d$ equation has support at most
\[
 \left(\frac1{p-1}+\varepsilon\right)d.
\]
More strongly, for every fixed real $\eta>0$, it remains complete
under the additional restriction
\[
 K=(p-1)\sigma-d\leq d^\eta,
\]
where $\sigma$ is the largest monomial support after discarding
monomials with an exponent at least $p$. The reductions have
polynomial bit size for fixed $p,\varepsilon,\eta$.
\end{theorem}
\begin{proof}
Start with the matrix in the proof of
Theorem~\ref{sp:thm:odd-linears}. Adjoin identity rows and columns
until its size $r$ is divisible by $p-1$. This adds at most $p-2$
rows and preserves the permanent. Keep the notation
\[
 P=4^{-t}\prod_{i=1}^rL_i,\qquad
 Z=\prod_{i=1}^r x_i,\qquad [Z]P=S.
\]
There are $p-1$ tag variables $t_1,\ldots,t_{p-1}$, and
$u=r+p-1$ variables before padding. Choose a power $Q$ of $p$ with
$Q\geq u+2$ and put
\[
 m=Q-1,\qquad d=Q+1,\qquad
 b=\frac{m-u}{p-1}.
\]
The integer $b$ is nonnegative, since $r$ and $Q-1$ are divisible
by $p-1$. Partition the $m-u$ new coordinates into $p-1$ blocks
$B_1,\ldots,B_{p-1}$, each of size $b$, and write
$W_j=\prod_{v\in B_j}v$. All $m$ coordinates, including the original
variables, tags and blocks, are denoted collectively by $x$.
Define the degree-$m$ form
\[
 A=-t_1^{p-1}P W_1^{p-1}
       +\sum_{j=2}^{p-1}t_j^{p-1}Z W_j^{p-1}+x_1^m
\]
and the degree-$d$ form in $d$ variables
\[
 F=yzA+y^{Q+1}+\sum_{i=1}^m x_i^{Q+1}-x_1z^Q.
\]
The tag and block coordinates are distinct from $x_1$.

We check smoothness directly because the padding has changed.
As $p\mid Q$, the identity
$yF_y-zF_z=y^{Q+1}$ forces $y=0$ at a common zero of the
derivatives. The derivatives in $x_i$ then give $x_i=0$ for $i\geq2$
and $x_1^Q=z^Q$. A projective common zero would consequently have
$x_1z\ne0$. Every tagged summand of $A$ vanishes on this coordinate
axis, whereas $A=x_1^m$ there. Thus $F_y=zA\ne0$, a contradiction.
The hypersurface is smooth; positive-dimensional smooth hypersurfaces
in projective space of this dimension are geometrically integral,
as in Theorem~\ref{thm:stabilizer}.

The pure powers and $x_1z^Q$ cannot contribute to the Hasse coefficient,
so that coefficient is $[(x_1\cdots x_m)^{p-1}]A^{p-1}$.
The term $x_1^m$ cannot contribute either, since $m\geq p$.
Every tag must occur exactly once. Each block variable then acquires
exponent $p-1$ from exactly its own summand, and each original
variable must be supplied once by $P$, since the other $p-2$
summands supply $Z^{p-2}$. The coefficient is therefore
\[
 -(p-1)![Z]P=S.
\]

Replace $\prod_iL_i$ by $\prod_iw_i$ and impose the same independent
three-term partial-sum equations as before. This produces a smooth
complete intersection of multidegree $(1^k,d)$ with the desired
scalar. The nonlinear equation still has at most $d+p$ monomials.
Each tagged monomial now has support at most
\[
 \sigma=r+b+3,
\]
counting the $r$ original or row variables, one tag, $b$ block
variables, and $y,z$. All other monomials have support at most three.
The tagged monomials have individual exponents at most $p-1$;
the other terms are removed by the deletion rule in
Theorem~\ref{sp:thm:support-strip}. Consequently
\[
 (p-1)\sigma-d=(p-2)(r+2),
 \qquad
 \sigma-\frac d{p-1}=\frac{(p-2)(r+2)}{p-1}.
\]
These quantities do not grow with $b$.

For the first assertion choose $Q$ so that
$d\geq(p-2)(r+2)/((p-1)\varepsilon)$ in addition to the earlier
conditions. The least suitable power of $p$ is $O_{p,\varepsilon}(r)$.
For the second, enlarge $Q$ until
$d^\eta\geq(p-2)(r+2)$ as well. This still gives $d$ polynomial
in $r$ for fixed $\eta>0$; equivalently one can choose a fixed
integer padding exponent larger than $1/\eta$.
The $O(d)$ pure powers and the block exponent vectors thus have
polynomial total bit length. The same normalized-point verifier
proves membership, completing both assertions.
\end{proof}

\section{Effective projection and the remaining presentation questions}
\label{sp:sec:compression}
An effective projection must specify its field, its center, and equations
for the image. The existence of a general center addresses only the
geometric part of that task.

\begin{proposition}[A field bound for an embedding projection]
\label{sp:prop:projection-field}
Let the hypersurface of Theorem~\ref{sp:thm:reembedding} be defined over
$\mathbb F_q$, and let an embedding by $aH$, $a\geq2$, have image in
$\PP^R$. Put $E=a^n\delta$. If $R>2n+1$, an embedding projection to
$\PP^{2n+1}$ exists over every extension $\mathbb F_{q^e}$ such that
\[
 q^e>(2n+2)E^2.
\]
One may therefore take
\[
 e=\left\lceil\log_q\bigl((2n+2)E^2+1\bigr)\right\rceil.
\]
When $R=2n+1$, the identity already suffices. For a presentation of
Theorem~\ref{sp:thm:trinomials}, $a=p^h$, so this extension degree is
$O(nh+n\log_p\delta+\log_p(n+2))$ over $\Fp$.
\end{proposition}
\begin{proof}
Let $S$ be the closed secant variety of the embedded $X$. We proved
$\dim S=2n+1$. Its degree $\Delta$ is at most $E^2$.
To see the bound, put two copies of $X$ in disjoint coordinate subspaces
of $\PP^{2R+1}$. Their join has coordinate ring the graded tensor
product of the two coordinate rings. The Hilbert series multiply, so
its dimension is $2n+1$ and its degree is $E^2$. The linear rational
projection $[u:v]\mapsto[u+v]$ has image $S$ and finite generic fibre.
Linear projection cannot increase degree: the isolated inverse images
of a general complementary linear section have total degree at most
the degree of the source, by the linear-section form of B\'ezout's
inequality. Thus $\Delta\leq E^2$.

Write $s=2n+1$. The Chow equation of $S$ in $s+1$ blocks of hyperplane
coefficients is nonzero and has degree $\Delta$ in each block. Its
vanishing means that the corresponding hyperplanes have a common
point on $S$. For completeness, the incidence of a point of $S$ and
$s+1$ hyperplanes through it is irreducible and has image of codimension
one in the product of dual projective spaces. Fixing $s$ general
hyperplanes leaves $\Delta$ points, and the last block's divisor is
the union of the $\Delta$ hyperplanes dual to those points. This gives
the stated degree in every block and defines the Chow equation over
the field of definition of $S$.

Its total degree is $(s+1)\Delta\leq(2n+2)E^2$. A nonzero polynomial
of total degree less than the size of a finite field cannot vanish
on the whole affine grid; induction in the number of variables proves
this from the univariate root bound. Thus over the indicated extension
there is a coefficient matrix on which the Chow equation is nonzero.
Its rows are independent, since otherwise their common linear space
would meet $S$ by the projective dimension theorem. Their common zero
is a center of codimension $s+1$ missing $S$. Separation of length-two
subschemes makes the projection a closed immersion. The expression for
$e$ and the last assertion follow by taking logarithms.
\end{proof}

\begin{corollary}[Randomized selection of a projection]
\label{sp:cor:random-projection}
In Proposition~\ref{sp:prop:projection-field}, suppose $R>2n+1$ and
an explicit extension of cardinality $Q=q^e$ is available. A uniformly
sampled $(2n+2)\times(R+1)$ matrix defines a closed-immersion projection
with probability at least
\[
 1-\frac{(2n+2)E^2}{Q}.
\]
Thus $Q>2(2n+2)E^2$ gives success probability greater than one half.
Exact uniform sampling takes expected polynomial bit time in
$n,R,\log Q$, once the embedding
and field are explicitly represented.
\end{corollary}
\begin{proof}
The nonzero Chow obstruction polynomial from the preceding proof has
total degree at most $(2n+2)E^2$. For a polynomial of total degree
$D$ in $M$ variables, induction using the univariate root bound gives
at most $D Q^{M-1}$ zeros on $\mathbb F_Q^M$. Applying this to the
matrix entries proves the probability estimate. On every nonzero
value its center avoids the closed secant variety and the projection
is a closed immersion, as proved above.
For a bounded-bit variant, use the coordinate sampler from
Corollary~\ref{sp:cor:unspecified-cofactor}. Its maximum point mass
is $2/Q$, so the failure bound becomes $2(2n+2)E^2/Q$.
Thus $Q>4(2n+2)E^2$ gives success probability greater than one half
with a bounded polynomial number of random bits.
\end{proof}

This is a Monte Carlo guarantee: the sampling procedure does not
certify its selected matrix or compute equations of the image. An
extension can also be constructed by random search for a monic
irreducible polynomial of degree $e$ over $\mathbb F_q$ and standard
finite-field irreducibility tests. The proportion of monic degree-$e$
polynomials that are irreducible is at least $1/(2e)$ for $e\geq2$,
by the usual count of elements of degree $e$; each test has polynomial
bit cost in $e,\log q$. This supplies expected polynomial extension
construction. None of these assertions supplies a center over the
original small field.

This proof is effective with large costs. Once a Chow equation is
computed, exhaustive search takes at most
$(q^e)^{(2n+2)(R+1)}$ matrix evaluations. Each evaluation or a direct
nonintersection test is exact. Computing the secant ideal and its Chow
equation is elimination work, and can itself be much larger than the
original circuit. No polynomial-time construction of a center over
the original small field follows from the bound.

There is also an explicit, though very large, bound for producing
scheme-theoretic image equations once a center has been chosen.
Suppose the source is given by homogeneous equations $J$ of degrees
at most $D$, put $s=2n+1$, and introduce target coordinates
$z_0,\ldots,z_s$. If $A$ is the projection matrix, use the equations
\[
 J(x),\qquad z-Ax=0,\qquad
 1-\sum_{i=0}^s t_i z_i=0.
\]
Eliminating $x$ and $t$ gives the saturated homogeneous ideal of the
image in the $z$ coordinates. The last equation removes the target
origin; contraction is saturation by $(z_0,\ldots,z_s)$. Since the
projection is a closed immersion on the projective source, this is
scheme-theoretic image formation, not replacement by a radical.

Here is a quantitative finite procedure. Homogenize these equations
with one additional variable. The number of variables is
$v=R+2s+4$ and their degrees are at most $D_0=\max(D,2)$.
Dub\'e's bound \citep[Theorem~8.2]{sp-dube} permits a Gr\"obner basis
for an elimination order of degree at most
\[
 B_0=\left\lceil2\bigl(D_0^2/2+D_0\bigr)^{2^{v-1}}\right\rceil.
\]
The vector spaces of homogeneous ideal elements through this degree
can be built from monomial multiples of the input generators.
Row reduction, with an elimination ordering on monomials of each
degree, gives a Gr\"obner basis because the bound ensures that every
minimal leading monomial occurs by this stage. Dehomogenization then
gives the required elimination ideal. If $M=\binom{v+B_0}{v}$ and
$b$ is the number of input equations in this construction, a crude
bound is $O(bM^3)$ field operations, with polynomial overhead for
monomial indexing. Homogeneous components and row reduction in the
target variables produce at most $\binom{s+1+B_0}{s+1}$ equations,
each of degree at most $B_0$ and with at most $\binom{s+B_0}{s}$ terms.
These bounds deliberately display the growth rather than treating
elimination as a unit-cost operation. They supply an effective
scheme-preserving procedure after the field extension and center have
been fixed, but not a polynomial-size sparse compression.

\paragraph{Comparison with a radical compression theorem.}
\citet[Theorem~3.2]{sp-crw} consider $m$ homogeneous degree-$D$ forms
in $r$ variables over $\mathbb F_q$. If $q>rD^{r-1}$, they construct
$r$ degree-$D$ forms with the same radical. For $D\geq2$ their time
bound is $m(rD^r)^{O(r)}\log^2q$; for $D=1$ it is
$O(\max(m,r)^3\log^2q)$. Their construction uses linear combinations
of the original forms, so it does not multiply their supports together.
To equalize different positive degrees without altering the projective
scheme, one may replace $f_i$ by all $x_j^{D-\deg f_i}f_i$:
on the chart $x_j\ne0$ the multiplier is a unit. This uses at most
$r$ times as many equations and preserves their term counts.
The resulting radical compression preserves rational points over every
extension, but does not certify equality of the scheme structures.
It also yields $r$ equations, not a fixed number, and its degree and
field bounds are not polynomial sparse-bit bounds when dimensions grow.

\paragraph{The questions left by these algorithms.}
For a fixed prime and fixed codimension, a long sparse list known only
to generate $I_X$ after saturation is not covered by
Theorem~\ref{sp:thm:generators}. Theorem~\ref{sp:thm:bounded-cofactors}
settles hypersurfaces when the unknown cofactor total degrees are bounded.
Theorem~\ref{sp:thm:one-cofactor} requires a fixed total-degree bound
for only one residual cofactor, and
Corollary~\ref{sp:cor:individual-scalar} covers bounded individual
cofactor degrees in polynomial time. More generally,
Theorem~\ref{sp:thm:sparse-cofactor} handles a designated residual
cofactor with polynomial numerical degree and a supplied polynomial
support bound. Corollary~\ref{sp:cor:unspecified-cofactor} removes
the designation in a Monte Carlo version under the corresponding
existence promise. For unrestricted residual degrees or supports,
the remaining problem is to compute the
coherent coefficient without enumerating all cofactor monomials or
a dense saturated ideal. In higher codimension,
Theorem~\ref{sci:thm:random-excess} gives a randomized algorithm
whenever the actual input ideal, after a finite field extension,
contains a homogeneous regular sequence of length $c$ and degree sum
at most $N+1+E$, with $E$ fixed. The sequence and its degrees need not
be supplied, and the other input degrees may be arbitrary binary
integers. The algorithm returns a correct scalar or reports failure;
repetition gives expected polynomial time. Theorems~\ref{sci:thm:gap}
and~\ref{sci:thm:balanced} retain deterministic conclusions under
their respective degree hypotheses. The remaining general saturation
gate allows unbounded least complete-intersection excess in the
actual input ideal; deterministic selection and identity testing are
additional questions in the randomized range. Replacing the actual
ideal by its saturation would remove an essential hypothesis.

For the same high-dimensional hypersurface, bounded-codimension
reembeddings must eventually use its original hyperplane class, by
Theorem~\ref{sp:thm:reembedding}. The questions are then whether a
nonsingular change of its original coordinates can make its equation
polynomially sparse, or whether polynomially many sparse multiples
can define it after saturation. Uniform construction of a single
sparse equation, or a bounded list, is excluded under the assumption
in Corollary~\ref{sp:cor:compression}. A long list is not excluded by
that argument. Adding gate coordinates gives the explicit answer in
Theorem~\ref{sp:thm:trinomials}, at the cost of growing codimension;
it changes the question.

The general coordinate-sparsification decision problem is already
hard: \citet[Theorem~1]{sp-bdss} prove NP-hardness over every field,
even when the input polynomial is homogeneous and explicitly sparse.
Their Theorem~3 proves hardness of equivalence to a polynomial of
monomial support at most $\sigma$, for fixed $\sigma\geq5$, under
the additional hypothesis that the characteristic is zero or greater
than $\sigma+1$. These results do not impose smoothness or the
anticanonical degree condition, so they do not settle the narrower
coordinate question above.

\citet[Theorems~1.2 and~1.5]{sp-bs-orbits} give polynomial-size
hitting sets and explicit lower bounds for sparse affine orbits in
characteristic zero. Their initial-form lemma, Lemma~1.6, does hold
in arbitrary characteristic. It bounds the Hasse-derivative dimension
of the lowest nonzero homogeneous component after most shifts along
a rational normal curve. That particular lower-bound mechanism has
a direct limitation for a smooth homogeneous hypersurface equation
$g$: at every nonzero affine point $a$, the initial form of $g(a+x)$
is either a nonzero constant or a nonzero linear form. Indeed the
first case is $g(a)\ne0$, and projective smoothness gives the second
when $g(a)=0$. Its derivative space therefore has dimension at most
two. Only the cone vertex can have a higher-order initial form.
Their hard examples instead use powers of a cubic singular along
the rational normal curve. Thus neither the field hypothesis nor
that geometric mechanism supplies the desired lower bound for the
smooth fixed-characteristic family here.

A different smooth projective scheme with the same scalar must specify
which coherent group carries it and how that scalar is read from all
its equations. For anticanonical complete intersections with an ideal-
generating sparse list, Theorem~\ref{sp:thm:generators} rules out a
uniform hard family at fixed $p,c$ under the same complexity assumption.
For bounded equations the point-count obstruction applies more broadly.
For long lists requiring saturation, or for other coherent groups with
no proved point-count readout, these arguments give no conclusion.
Theorem~\ref{sp:thm:support} provides a tractable range of logarithmic
monomial support for arbitrary zero schemes with many linear equations.
For anticanonical complete intersections,
Theorem~\ref{sp:thm:support-strip} also handles support within
$O(\log L)$ of the minimum compatible with a nonzero scalar.
Theorem~\ref{sp:thm:support-hardness} proves hardness even with excess
$K\leq d^\eta$ for each fixed $\eta>0$. The remaining parameter gap
is between logarithmic excess, which is tractable, and arbitrarily
small fixed powers of the dimension, where hardness persists.
These statements do not classify every intermediate growth rate.
